\section{CaliberGraph and CaliberGraph-R}\label{sec:method}

CaliberGraph takes $(\mathcal S,\Sigma)$, $Q_g$, $Q_p$, a target tier $\mathcal A$ and a wall-clock budget $B$, and returns a certified witness or nothing (Fig.~\ref{fig:overview}). Every candidate database is executed by SQLite itself, so the method inherits SQLite's type affinity, three-valued logic, string and date functions, integer division and tie-breaking instead of modelling them. Four components guide the search; each can be disabled for ablation.

\begin{figure}[t]
\centering
\begin{tikzpicture}[font=\scriptsize,>=Latex,
  box/.style={draw,rounded corners=2pt,align=center,inner sep=2pt,minimum height=8mm,text width=33mm},
  io/.style={box,fill=black!4},
  cert/.style={box,very thick}]
\node[io] (in) at (0,0) {DDL $\mathcal S,\Sigma$; gold $Q_g$, prediction $Q_p$};
\node[box,right=3mm of in] (an) {\textbf{Pair analysis}: resolved columns, equalities, atoms, diff anchors};
\node[box,below=4mm of an] (jg) {\textbf{Join-graph completion}: union--find classes; probe = one joined tuple; \textbf{SQLite value domain}};
\node[box,left=3mm of jg] (iso) {\textbf{Isolation \& growth}: single probes; duplicate, tie, decoy, fan-out, window};
\node[box,below=4mm of iso] (ev) {\textbf{Search loop}: SQLite evaluation in a savepoint; mutation};
\node[box,dashed,right=3mm of ev] (rs) {\textbf{Real sub-databases} (CaliberGraph-R): anchored retrieval, join completion, real growth from $D_0$, $\Pi$};
\node[cert,below=5mm of ev] (ce) {\textbf{Independent certifier}: original DDL + SQL; contracts; tiers $\mathcal A_0$--$\mathcal A_2$, $\mathcal A_R$};
\node[io,right=3mm of ce] (out) {certified witness, or none within budget};
\draw[->] (in)--(an);
\draw[->] (an)--(jg);
\draw[->] (jg)--(iso);
\draw[->] (iso)--(ev);
\draw[->,dashed] (rs)--(ev);
\draw[->] (ev)-- node[right,font=\scriptsize]{results differ} (ce);
\draw[->] (ce)-- node[below,font=\scriptsize]{acc.} (out);
\draw[->,dashed] (ce.west) -- ++(-2mm,0) |- node[pos=0.25,left,font=\scriptsize,align=right]{rej.} (ev.west);
\end{tikzpicture}
\caption{CaliberGraph pipeline. Candidates built from join-graph probes are evaluated by SQLite; a candidate that separates the queries is returned only if the independent certifier accepts it at the target tier. CaliberGraph-R (dashed) first searches real sub-databases of the benchmark database.}
\label{fig:overview}
\end{figure}


\textbf{Pair analysis and diff anchors.} Both queries are parsed in the SQLite dialect; column references are resolved to $(T,c)$ through aliases, subqueries and CTEs. For each query $Q$ we collect its tables $\mathrm{tab}(Q)$, its column equalities $E(Q)$ (including \texttt{IN}-subquery pairs) and its \emph{atoms} $\mathit{At}(Q)$: comparisons, \texttt{LIKE}, \texttt{IN}, \texttt{BETWEEN} and \texttt{IS [NOT] NULL} between a column, optionally wrapped in one function (\texttt{strftime}, \texttt{substr}, \texttt{lower}, \texttt{CAST}, $\pm k$, \dots), and literals. The \emph{anchor atoms} are those whose normal form (resolved names, canonical operator) occurs in only one query, and the \emph{anchor columns} add every column referenced, projected, ordered, grouped or aggregated in only one query. Anchors localise where the queries can disagree.

\textbf{Join-graph completion.} Let $\mathcal T^{*}$ be $\mathrm{tab}(Q_g)\cup\mathrm{tab}(Q_p)$ closed under foreign-key parents. Columns of $\mathcal T^{*}$ are partitioned into classes by the closure of $E(Q_g)\cup E(Q_p)$ and of the foreign keys, computed with union--find. A \emph{probe} assigns one value per class; its materialisation contains one row per referenced table, so that all equated columns agree, and a repair step adds parent rows for dangling foreign keys, fills \texttt{NOT NULL} columns and drops duplicate keys. Key classes receive distinct values in distinct probes.

\begin{proposition}[Reachability]
If a probe makes every atom of the outermost \texttt{WHERE} clause of $Q$ true, its rows form a joined tuple of $Q$ that satisfies the join conditions and filters of $Q$ in the repaired database.
\end{proposition}
\noindent Each equality of $E(Q)$ relates columns of one class, which take one value; repair only adds rows, fills \texttt{NULL}s in \texttt{NOT NULL} columns and drops later duplicates, so the probe's rows are unchanged. Random rows satisfy a join only by chance, which is why fuzzing draws keys from tiny domains; probes reach the output by construction whenever their filters hold.

\textbf{SQLite-aware value domain.} Each class receives a pool of values from the atoms on its columns: each literal $v$ and $v\pm1$, $v\pm0.5$ for numeric columns (integer division, \texttt{CAST}); case and padding variants of strings (\texttt{=} is case-sensitive, \texttt{LIKE} is not); pre-images under wrapping functions (e.g.\ \texttt{'2012-06-15 10:00:00'} for \texttt{strftime('\%Y',$c$)='2012'}); matching and non-matching \texttt{LIKE} instances; and values strictly between two comparable literals. Values are coerced to the declared type, and text of an unobserved shape is replaced by an observed sample.

\textbf{Isolation and controlled growth.} Aggregates, \texttt{DISTINCT}, \texttt{GROUP BY} and \texttt{ORDER BY}\dots\texttt{LIMIT} can \emph{mask} a difference. The search therefore starts from single-probe databases, in which one joined tuple determines the output, and grows them only with operators that target a masking mechanism (Table~\ref{tab:operators}).

\begin{table}[t]
\centering\footnotesize
\caption{Growth operators (each followed by repair).}\label{tab:operators}
\begin{tabularx}{\linewidth}{@{}l>{\raggedright\arraybackslash}X@{}}
\toprule
Operator & Construction: difference exposed\\
\midrule
duplicate & row copy with fresh keys: \texttt{DISTINCT}, \texttt{COUNT(DISTINCT)}\\
fan-out & $n$ child rows of a join: multiplicity, \texttt{HAVING COUNT}$>n$\\
tie probe & fresh copy tied on ordered columns, both insertion orders: \texttt{LIMIT} tie-breaking\\
window & $\ell{+}o{+}1$ tie probes: rows at \texttt{LIMIT}~$\ell$ \texttt{OFFSET}~$o$ borders\\
literal cover & one probe per equality literal: \texttt{CASE} branches, \texttt{IN} lists\\
decoy & a row that joins nothing: outer joins, anti-joins, \texttt{NOT IN}\\
violating & a probe failing the filters: empty inputs, \texttt{MAX} over no row\\
\bottomrule
\end{tabularx}
\end{table}

\begin{algorithm}[t]
\caption{CaliberGraph (CaliberGraph-R adds line 4)}\label{alg:main}
\footnotesize
\begin{algorithmic}[1]
\Require DDL $(\mathcal S,\Sigma)$, queries $Q_g,Q_p$, tier $\mathcal A$, budget $B$; for $\mathcal A_R$: $D_0$, $\Pi$
\Ensure witness $W\in\mathcal A$ with $\neg\mathcal C_{\mathrm{set}}(Q_g(W),Q_p(W))$, or $\bot$
\State analyse $Q_g,Q_p$: equalities $E$, atoms $\mathit{At}$, anchors
\State union--find column classes over $\mathcal T^{*}$; a value pool per class
\State $\mathit{Pool}\gets\emptyset$
\State \textbf{if} $\mathcal A=\mathcal A_R$: put real sub-databases (anchored retrieval, join completion, real growth) first in the stream
\For{each $D$ in the candidate stream \textbf{while} time $<B$}
  \State $(r_g,r_p)\gets$ run $Q_g,Q_p$ on $D$ in a savepoint
  \If{$\mathrm{set}(r_g)\neq\mathrm{set}(r_p)$ \textbf{and} \Call{Certify}{$D,\mathcal A$}}
    \State \Return $D$ \Comment{re-certified on a fresh connection}
  \EndIf
  \State add $D$ to $\mathit{Pool}$ with weight $=$ number of non-empty results among $r_g,r_p$
\EndFor
\State \Return $\bot$
\Statex \emph{Candidate stream, in order:} single probes for $Q_g$, $Q_p$ and both; two-probe combinations; growth operators (Table~\ref{tab:operators}); boundary and \texttt{NULL} variants of anchor classes; then mutations of parents drawn from $\mathit{Pool}$ by weight, each followed by repair
\end{algorithmic}
\end{algorithm}

\textbf{Search and certification.} The candidate stream (Algorithm~\ref{alg:main}) enumerates single-probe databases satisfying $Q_g$, $Q_p$ and both, then two-probe combinations, the growth structures and boundary and \texttt{NULL} variants of anchor classes, and finally runs a mutation loop until the budget expires (parents weighted by how many queries return rows; 70\,\% of value changes on anchor columns). Each candidate is evaluated in a savepoint of one in-memory SQLite connection (a few hundred microseconds). A candidate on which the results differ is returned only if the certifier accepts it at the target tier, and the certifier is re-run on a fresh connection before reporting. \emph{Soundness} is immediate: a returned $W$ is in $\mathcal A$ and separates the original strings on SQLite. The guarantee is one-sided and relative to the correctness of SQLite and the certifier. CaliberGraph+SC (``SC'': SpotIt+ constraints) additionally uses the constraints mined by SpotIt+ for value pools and acceptance, so that it has the same knowledge as SpotIt+; no other self-run tool uses them.

\subsection{Data-grounded search for $\mathcal A_R$}\label{sec:realsearch}
CaliberGraph-R changes where candidate rows come from. It reads $D_0$ (through an indexed read-only copy) and $\Pi$, and searches \emph{real sub-databases} before synthesis. A sub-database of $D_0$ that is closed under the inclusion dependencies of $\Pi$ satisfies the profile, except for rare rows that violate an approximate cross-key dependency.
\begin{itemize}
\item \emph{Diff-anchored retrieval.} The atoms of each query on table $T$ are translated back to SQL predicates $\varphi_g^T,\varphi_p^T$; rows with $\varphi_g^T\wedge\neg\varphi_p^T$ or $\varphi_p^T\wedge\neg\varphi_g^T$ are retrieved first (the filters disagree on them), then rows satisfying both, one, or neither.
\item \emph{Join completion.} Each seed row is completed along the join graph with real related rows, preferably of the same predicate mode, and closed under foreign keys, undeclared foreign keys and total participation.
\item \emph{Real growth.} Sub-databases are combined (a row only one query selects next to rows both select) and grown with real rows that tie on an \texttt{ORDER BY} column, the first $k{+}1$ rows of a \texttt{LIMIT}~$k$ window, real \texttt{NULL}s and real duplicates; they are shrunk by single-row deletions, and an entity can be cloned with the rows that reference it.
\item \emph{Constrained synthesis.} If no real sub-database separates the queries, the synthetic search continues from the evaluated sub-databases with pools filtered by $\Pi$, untouched columns filled from real template rows, and a repair step that also enforces $\Pi$.
\end{itemize}

\textbf{Example (development pair).} For BIRD question 1020 (``Which player has the highest overall rating?''), $Q_g$ orders \texttt{Player\_Attributes} by \texttt{overall\_rating DESC LIMIT 1}, whereas the Alpha-SQL prediction selects every row equal to the maximum rating. Both return one player on $D_0$. CaliberGraph-R retrieves a real attribute record (rating 84) with its player, then adds a \emph{real} record of another player with the same rating. On this four-row database, every row of which occurs in BIRD, $Q_g$ returns one player and $Q_p$ two; the witness is certified at $\mathcal A_R$ after 28 candidates and 0.22\,s. Whether the prediction or the gold query is right depends on how the question treats ties, which is what attribution judges.
